% =============================================================
% Capitolo: User Profiling — Memoria a Lungo Termine dell'Agente
% Progetto: MCP_ECOM — Agente e-commerce basato su MCP
% =============================================================

\chapter{User Profiling e Memoria a Lungo Termine}
\label{chap:user-profiling}

% ─────────────────────────────────────────────────────────────
\section{Visione d'Insieme}
\label{sec:profiling-intro}
% ─────────────────────────────────────────────────────────────

Il sistema di \textbf{User Profiling} rappresenta la memoria a lungo termine (LTM) dell'agente. Invece di trattare ogni sessione come un evento isolato, l'agente apprende passivamente dalle interazioni dell'utente (ricerche, click, analisi) per costruire un profilo persistente. 

Questo profilo viene utilizzato per:
\begin{itemize}
    \item Raffinare il \textbf{Reranking} dei prodotti (\S\ref{sec:reranking}), dando priorità ai brand e alle condizioni preferite.
    \item Personalizzare il \textbf{Budgeting}, suggerendo prodotti in linea con la fascia di prezzo storica dell'utente per quella specifica categoria.
    \item Adattare lo \textbf{stile comunicativo} in base alla profondità di interazione rilevata.
\end{itemize}

Il servizio agisce in modo trasparente e non distruttivo: le preferenze esplicite impostate dall'utente (via MCP tool) coesistono con quelle apprese automaticamente, avendo sempre la priorità.

% ─────────────────────────────────────────────────────────────
\section{Architettura e Segnali di Apprendimento}
\label{sec:profiling-signals}
% ─────────────────────────────────────────────────────────────

Il modulo \texttt{user\_profiling.py} estrae segnali semantici dall'output del parser (\S\ref{sec:query-parsing}) e dai risultati della ricerca live. I segnali principali sono:

\begin{enumerate}
    \item \textbf{Price Signal}: Estrazione dei vincoli di prezzo per identificare la disponibilità di spesa.
    \item \textbf{Condition Signal}: Rilevamento della preferenza tra Nuovo, Usato o Ricondizionato.
    \item \textbf{Category Affinity}: Tracciamento delle categorie merceologiche più visitate.
    \item \textbf{Interaction Depth}: Monitoraggio del livello di "engagement" dell'utente.
\end{enumerate}

\subsection{Exponential Moving Average (EMA) per i Budget}
\label{subsec:profiling-ema}

Per evitare che una singola ricerca outlier (es. un utente che solitamente cerca smartphone economici ma curiosa su un iPhone 15 Pro) stravolga il profilo, il sistema utilizza una \textbf{Media Mobile Esponenziale (EMA)} per l'aggiornamento dei budget contestuali:

\begin{equation}
    Budget_{new} = Budget_{old} \times 0.80 + Signal \times 0.20
\end{equation}

Questo garantisce che il profilo sia resiliente ai "rumori" di breve termine ma sensibile ai cambiamenti di abitudine consolidati.

% ─────────────────────────────────────────────────────────────
\section{Modellizzazione delle Preferenze}
\label{sec:profiling-models}
% ─────────────────────────────────────────────────────────────

Le preferenze dell'utente sono salvate nel database in formati serializzati (JSON o Counter-Strings) per permettere una rappresentazione ricca e flessibile.

\subsection{Preferenze di Brand e Condizione}
\label{subsec:profiling-brands}

I brand non sono memorizzati come una semplice lista, ma come un dizionario con conteggio di "hit". Solo i brand con un'affinità rilevante vengono iniettati nel contesto del prompt per evitare allucinazioni o inquinamento della UI.

\begin{verbatim}
# Esempio di serializzazione nel DB
favorite_brands: "Apple:12,Samsung:4,Sony:1"
condition_preference: "new:8,refurbished:2"
\end{verbatim}

Il sistema di \textbf{maggioranza qualificata} richiede almeno 3 segnali passati e una soglia del 50\% di occorrenze prima di considerare una condizione (es. "Nuovo") come "Dominante" per il filtraggio automatico.

\subsection{Budget Contestuali (Contextual Budgets)}
\label{subsec:profiling-context-budget}

Una delle funzionalità più avanzate è il budget contestuale per categoria o brand. L'utente potrebbe avere un budget di 20€ per le "Cover" ma di 800€ per gli "Smartphone". Il sistema distingue tra:
\begin{itemize}
    \item \texttt{auto\_cat:}: Budget appreso automaticamente per la categoria.
    \item \texttt{auto\_brand:}: Budget appreso per uno specifico brand.
    \item \texttt{manual:}: Limiti impostati esplicitamente dall'utente che sovrascrvono gli "auto".
\end{itemize}

% ─────────────────────────────────────────────────────────────
\section{Interaction Depth e Engagement}
\label{sec:profiling-depth}
% ─────────────────────────────────────────────────────────────

L'agente classifica l'utente in tre tier di profondità di interazione (Tabella~\ref{tab:interaction-depth}). Questo tier determina quanto dettagliata debba essere la ricerca o se l'agente debba proporre proattivamente analisi di trust dei seller.

\begin{table}[htbp]
\centering
\caption{Tier di Interaction Depth e relativi trigger}
\label{tab:interaction-depth}
\begin{tabular}{lll}
\hline
\textbf{Tier} & \textbf{Descrizione} & \textbf{Trigger Principali} \\
\hline
\texttt{browser} & Navigazione superficiale & Solo ricerche generali. \\
\texttt{researcher} & Analisi comparativa & Visualizzazione dettagli, comparazione, analisi seller. \\
\texttt{power\_buyer} & Intento d'acquisto elevato & Contatto venditore, molteplici analisi profonde. \\
\hline
\end{tabular}
\end{table}

% ─────────────────────────────────────────────────────────────
\section{Iniezione nel Prompt (Context Building)}
\label{sec:profiling-context}
% ─────────────────────────────────────────────────────────────

Prima di ogni interazione con l'LLM, il metodo \texttt{build\_profile\_context} trasforma i dati grezzi del database in un oggetto strutturato (Snapshot) che viene inserito nel System Prompt.

\begin{verbatim}
{
  "favorite_brands": "Apple, Sony",
  "price_preference": "500",
  "condition_preference": "new",
  "category_affinities": ["Smartphone", "Audio"],
  "interaction_depth": "researcher"
}
\end{verbatim}

Questo permette all'agente di formulare risposte del tipo: \textit{"Dato che solitamente preferisci prodotti Apple e cerchi condizioni pari al nuovo, ho selezionato questo iPhone ricondizionato di Grado A++ che rientra perfettamente nel tuo budget abituale."}


% ─────────────────────────────────────────────────────────────
\section{Collegamento dell'Account e Identità Persistente}
\label{sec:profiling-account}
% ─────────────────────────────────────────────────────────────

Il sistema non si limita a un database isolato, ma crea un ponte tra l'identità dell'utente nell'applicazione e la sua sessione reale su eBay. Questo avviene attraverso il meccanismo del \textbf{session\_id}, che funge da chiave primaria per la risoluzione dell'utente nel database tramite la funzione \texttt{resolve\_user\_by\_id}.

\subsection{Sincronizzazione tra Database e Sessione Browser}
\label{subsec:profiling-sync}

A differenza dei sistemi e-commerce tradizionali, il progetto MCP\_ECOM gestisce un'identità "fluida":
\begin{itemize}
    \item \textbf{Livello Database}: Vengono memorizzate le preferenze analitiche (brand, budget storici, interaction depth).
    \item \textbf{Livello Browser (Playwright)}: Grazie alla connessione via protocollo CDP o all'uso di profili persistenti (\S\ref{subsec:pw-contact-conn}), l'agente opera con i reali cookie di sessione dell'utente.
\end{itemize}

Questo collegamento permette all'agente di conoscere, ad esempio, se un utente ha già salvato un oggetto nella sua watchlist reale o di inviare messaggi a nome dell'utente senza mai visualizzare le sue credenziali, garantendo la massima privacy e sicurezza.

% ─────────────────────────────────────────────────────────────
\section{Personalizzazione Esplicita: Il Tool update\_user\_preferences}
\label{sec:profiling-manual}
% ─────────────────────────────────────────────────────────────

Oltre all'apprendimento passivo descritto in precedenza, il sistema offre all'utente la possibilità di "istruire" direttamente l'agente tramite il tool \texttt{update\_user\_preferences}. Questo strumento consente una personalizzazione granulare che ha la priorità assoluta sugli algoritmi di profiling automatico.

\subsection{Override dei Brand e delle Condizioni}
\label{subsec:profiling-brand-custom}

L'utente può dichiarare esplicitamente le proprie preferenze (es. \textit{"Compro solo scarpe Nike"}). In questo caso, l'agente invoca il tool impostando il campo \texttt{favorite\_brands}. 
A livello di database, questa azione:
\begin{enumerate}
    \item Sovrascrive il contatore di affinità generato dalle ricerche precedenti.
    \item Inserisce il brand specificato come segnale dominante nel prompt di sistema.
    \item Forza il modulo di Reranking (\S\ref{sec:reranking}) ad assegnare il punteggio massimo di rilevanza ai prodotti del brand scelto.
\end{enumerate}

Lo stesso principio si applica alla \texttt{condition\_preference}: se l'utente imposta manualmente "Usato", il sistema inietta nel DB un contatore fittizio ad alto valore (es. \texttt{used:99}) per silenziare l'apprendimento automatico e bloccare la preferenza su quel valore.

\subsection{Personalizzazione del Tono e delle Istruzioni}
\label{subsec:profiling-tone}

Il profilo utente include campi dedicati alla \emph{metacomunicazione}, permettendo di personalizzare l'esperienza d'uso oltre il semplice acquisto:
\begin{description}
    \item[conversation\_tone]: L'utente può scegliere tra toni \texttt{neutral}, \texttt{friendly}, \texttt{professional} o addirittura \texttt{zen} per una comunicazione più sintetica.
    \item[custom\_instructions]: Un campo di testo libero dove l'utente può inserire regole fisse (es. \textit{"Non suggerirmi mai prodotti che costano più di 500€"} o \textit{"Ignora i venditori con meno di 10 feedback"}).
\end{description}

\begin{figure}[htbp]
    \centering
    \fbox{\parbox{0.85\linewidth}{
        \textbf{Tool Invocato:} \texttt{update\_user\_preferences} \\
        \textbf{Argomenti:} \texttt{\{favorite\_brands: "Nike, Adidas", price\_preference: "Premium", conversation\_tone: "professional"\}} \\
        \textbf{Risultato:} Il profilo viene aggiornato e viene inviata una notifica di \texttt{resource\_updated} al client MCP per rinfrescare lo stato della UI.
    }}
    \caption{Esempio di flusso per la personalizzazione esplicita del profilo.}
    \label{fig:profiling-update-tool}
\end{figure}

% ─────────────────────────────────────────────────────────────
\section{Osservabilità: La Risorsa profile://}
\label{sec:profiling-resource}
% ─────────────────────────────────────────────────────────────

Per garantire la trasparenza, il sistema espone il profilo utente come una \textbf{Risorsa MCP} accessibile tramite l'URI template \texttt{profile://\{session\_id\}}. Ogni volta che le preferenze cambiano (sia per apprendimento automatico che per input manuale), l'agente riceve una notifica e può "rileggere" il profilo aggiornato. Questo chiude il loop di feedback tra azione, apprendimento e consapevolezza dell'AI, rendendo l'agente capace di giustificare le proprie scelte citando esplicitamente le preferenze salvate dell'utente.

